\begin{proofbox}
	\textbf{\textcolor{CProof}{思路提示}}

	从两点距离公式出发，利用椭圆方程消去 $y_0$，通过代数恒等变形得到完全平方式，再开方即得结论。

	\medskip
	\textbf{\textcolor{CProof}{正式推导}}
	\begin{description}[style=nextline, leftmargin=2.8em, labelsep=0.5em, itemsep=0.55em]
		\item[\textbf{步骤一（平方展开）：}] 写出 $|PF_2|$ 的平方：
			\[
				|PF_2|^2 = (x_0 - c)^2 + y_0^2.
			\]
			\par\textbf{理由：} 两点间距离公式 $|AB| = \sqrt{(x_1-x_2)^2+(y_1-y_2)^2}$。

		\item[\textbf{步骤二（代入消元）：}] 由椭圆方程 $\frac{x_0^2}{a^2}+\frac{y_0^2}{b^2}=1$ 得 $y_0^2 = b^2 - \frac{b^2}{a^2}x_0^2$，其中 $b^2 = a^2 - c^2$。
			\[
				|PF_2|^2 = (x_0-c)^2 + a^2 - c^2 - \frac{a^2-c^2}{a^2}x_0^2.
			\]
			\par\textbf{理由：} 将 $y_0^2$ 表达式代入，并将 $b^2$ 用 $a^2-c^2$ 替换。

		\item[\textbf{步骤三（合并化简）：}] 将上式展开并整理：
			\[
				\begin{aligned}
					|PF_2|^2 &= (x_0^2 - 2c x_0 + c^2) + a^2 - c^2 - \left(1-\frac{c^2}{a^2}\right)x_0^2 \\
					&= x_0^2 - 2c x_0 + a^2 - x_0^2 + \frac{c^2}{a^2}x_0^2 \\
					&= a^2 - 2c x_0 + \frac{c^2}{a^2}x_0^2.
			\end{aligned}
			\]
			\par\textbf{理由：} 合并同类项，注意 $b^2/a^2 = 1-c^2/a^2$。

		\item[\textbf{步骤四（配方开方）：}] 上式可化为完全平方：
			\[
				|PF_2|^2 = \left(a - \frac{c}{a} x_0\right)^2 = (a - e x_0)^2,
			\]
			因为 $e = c/a$。又因为 $a - e x_0 > 0$（$x_0 \le a$，$e<1$），所以
			\[
				|PF_2| = a - e x_0.
			\]
			\par\textbf{理由：} 完全平方公式；开方后取正值。

		\item[\textbf{步骤五（对称类推）：}] 同理计算 $|PF_1|^2 = (x_0 + c)^2 + y_0^2$，可得
			\[
				|PF_1|^2 = (a + e x_0)^2 \quad\Rightarrow\quad |PF_1| = a + e x_0.
			\]
			\par\textbf{理由：} 对称地将 $c$ 换为 $-c$，得到 $+2c x_0$ 项，配方后为 $(a + e x_0)^2$。

		\item[\textbf{步骤六（验证定义）：}] 检验和：
			\[
				|PF_1| + |PF_2| = (a + e x_0) + (a - e x_0) = 2a,
			\]
			与椭圆第一定义吻合，确认公式正确。
			\par\textbf{理由：} 椭圆定义 $|PF_1|+|PF_2|=2a$。
	\end{description}

	\medskip
	\textbf{\textcolor{CProof}{结论回扣}}

	\textbf{由此严格推导出椭圆的焦半径公式：$|PF_1| = a + e x_0$，$|PF_2| = a - e x_0$。} 该公式在解决椭圆距离问题时可大幅简化运算。
\end{proofbox}
